\documentclass{article}

% Language setting
% Replace `english' with e.g. `spanish' to change the document language
\usepackage[english]{babel}

% Set page size and margins
% Replace `letterpaper' with `a4paper' for UK/EU standard size
\usepackage[letterpaper,top=2cm,bottom=2cm,left=3cm,right=3cm,marginparwidth=1.75cm]{geometry}

% Useful packages
\usepackage{amsmath}
\usepackage{graphicx}
\usepackage[colorlinks=true, allcolors=blue]{hyperref}

\title{CSS}
\author{Dumitrita Sorohan}
\date{}

\begin{document}
\maketitle

\begin{abstract}

\end{abstract}

\section{Personal Contribution}

My contribution to the project focused on the implementation and validation of the process execution subsystem of the operating system scheduler simulator.  
More specifically, I worked on the development of the \texttt{ExecutionStage} and \texttt{Process} classes, together with the corresponding unit tests and runtime assertions required in the later project phases.

The implemented components are responsible for simulating the behaviour of user and system processes, handling execution stages, system calls, CPU affinity tracking, and process memory state management.


\section{Phase 1 -- Application Development}

\subsection{ExecutionStage Class}

The \texttt{ExecutionStage} class was designed to represent one execution segment of a user process.  
Each stage contains:

\begin{itemize}
    \item the number of CPU execution ticks;
    \item the duration of the following system call;
    \item or \texttt{None} if the stage is the final execution segment.
\end{itemize}

For example, the sequence:

\begin{verbatim}
(5 2 3 4 6)
\end{verbatim}

is interpreted as:

\begin{itemize}
    \item execute for 5 ticks, then perform a system call of 2 ticks;
    \item execute for 3 ticks, then perform a system call of 4 ticks;
    \item execute for 6 ticks and terminate.
\end{itemize}

The class also performs validation of the received values in order to prevent invalid execution stages from being created.



\subsection{Process Class}

The \texttt{Process} class represents both user processes and the special system process responsible for executing system calls.

The implementation supports:

\begin{itemize}
    \item parsing execution sequences into execution stages;
    \item execution progress through CPU ticks;
    \item transitions between execution stages;
    \item system call management;
    \item CPU affinity tracking;
    \item memory residency simulation.
\end{itemize}

\subsubsection{Execution Sequence Parsing}

One important part of the implementation was transforming the flat execution sequence into structured execution stages.

For example:

\begin{verbatim}
[5, 2, 3, 4, 6]
\end{verbatim}

is internally converted into:

\begin{center}
\begin{tabular}{|c|c|c|}
\hline
Stage & CPU Burst & System Call \\
\hline
1 & 5 & 2 \\
2 & 3 & 4 \\
3 & 6 & None \\
\hline
\end{tabular}
\end{center}

This functionality was implemented inside the private method:

\begin{verbatim}
_parse_sequence()
\end{verbatim}



\subsubsection{Process Execution Logic}

The execution logic was implemented through methods such as:

\begin{verbatim}
tick()
advance()
is_done()
is_waiting_for_sys_call()
\end{verbatim}

The \texttt{tick()} method advances the execution of the process by one unit of time, while \texttt{advance()} moves the process to the next execution stage after a system call finishes.

The implementation distinguishes clearly between user processes and the system process, since their runtime behaviour is different.



\subsubsection{System Call Queue}

For the system process, I implemented a queue-based mechanism for handling system calls requested by user processes.

System calls are stored using the helper class:

\begin{verbatim}
SysCallRequest
\end{verbatim}

The queue allows the system process to execute requests sequentially and resume the blocked user processes after completion.



\subsubsection{CPU Affinity Tracking}

Another implemented feature was CPU affinity tracking.  
Since the scheduler specification states that a process should preferably continue execution on the processor it used most recently, each process stores a history of CPU assignments.

The following methods were implemented:

\begin{verbatim}
record_cpu()
get_last_cpu()
get_cpu_history()
\end{verbatim}



\subsubsection{Memory State Management}

The process implementation also contains support for memory residency simulation.

The methods:

\begin{verbatim}
load_into_memory()
evict_from_memory()
is_in_memory()
\end{verbatim}

allow the scheduler and memory manager to determine whether a process is currently loaded into RAM or stored on disk.



\section{Phase 2 -- Unit Testing}

For testing, I used Python's built-in \texttt{unittest} framework.

The goal of the tests was to verify:

\begin{itemize}
    \item correct parsing of execution sequences;
    \item correct execution behaviour;
    \item proper stage transitions;
    \item correct system call handling;
    \item CPU history functionality;
    \item memory state transitions;
    \item handling of invalid inputs and edge cases.
\end{itemize}

\subsection{ExecutionStage Tests}

The tests for \texttt{ExecutionStage} verify:

\begin{itemize}
    \item valid object creation;
    \item invalid execution durations;
    \item invalid system call durations;
    \item correct handling of final stages;
    \item string representation behaviour.
\end{itemize}

Example:

\begin{verbatim}
with self.assertRaises(ValueError):
    ExecutionStage(0, 2)
\end{verbatim}



\subsection{Process Tests}

The \texttt{Process} class tests cover both user and system process behaviour.

The tested scenarios include:

\begin{itemize}
    \item execution sequence parsing;
    \item decrementing execution ticks;
    \item process completion;
    \item waiting for system calls;
    \item queueing multiple system calls;
    \item triggering user process advancement after system call completion;
    \item CPU history recording;
    \item memory loading and eviction.
\end{itemize}

Several edge cases were also tested, such as invalid stage advancement or executing a system process with an empty syscall queue.



\section{Phase 3 -- Assertions}

Assertions were added directly into the implementation code in order to validate program correctness during runtime.

The assertions were mainly used for:

\begin{itemize}
    \item preconditions;
    \item postconditions;
    \item class invariants;
    \item state consistency checks.
\end{itemize}

\subsection{Assertions in ExecutionStage}

The constructor validates that execution durations are always positive.

Example:

\begin{verbatim}
assert run_ticks > 0
assert sys_call_ticks is None or sys_call_ticks > 0
\end{verbatim}



\subsection{Assertions in Process}

A dedicated method named:

\begin{verbatim}
_check_invariants()
\end{verbatim}

was implemented to verify the internal consistency of a process object.

The invariants ensure that:

\begin{itemize}
    \item process identifiers are valid;
    \item execution stages remain consistent;
    \item runtime values never become negative;
    \item stage indexes remain inside valid bounds.
\end{itemize}

Example invariant:

\begin{verbatim}
assert 0 <= self.left_to_run <= current.run_ticks
\end{verbatim}



Additional assertions were inserted into methods such as:

\begin{verbatim}
tick()
advance()
add_sys_call()
record_cpu()
\end{verbatim}

These checks helped identify invalid runtime states during development and debugging.



\section{Conclusion}

Through this contribution, I implemented the core process execution functionality required by the operating system scheduling simulator.

The developed subsystem provides:

\begin{itemize}
    \item process execution modelling;
    \item system call handling;
    \item execution stage management;
    \item CPU affinity tracking;
    \item memory state simulation;
    \item runtime correctness validation using assertions;
    \item extensive unit testing coverage.
\end{itemize}

The implementation was intentionally kept modular and relatively simple in order to facilitate integration with the scheduler, memory manager, and graphical interface developed in the other parts of the project.


\end{document}
